\chapter{Literature Review}

Valerian steel, a rare and highly coveted material in the fantastical realms of metallurgy, has sparked the curiosity of scholars and blacksmiths alike for centuries. Despite its legendary status, the scientific community remains baffled by the peculiar properties of this metal. One of the earliest known studies, "A Brief History of Glowing Swords" by \cite{smithson1427}, suggests that valerian steel was primarily used for weaponry, though its precise composition remains a mystery \cite{smithson1427}.

Modern scholars, however, have expanded their focus beyond mere swordsmithing to explore valerian steel’s potential as an energy source. \cite{fizzlewrench2008} proposed that, when exposed to a full moon, valerian steel may generate a brief but intense burst of magical energy, capable of powering entire cities for a split second \cite{fizzlewrench2008}. This breakthrough, unfortunately, had no practical applications, as the magical surge typically destroyed any city it powered.

On the more serious side of research, \cite{thorne2012} argued that valerian steel’s resistance to both rust and magic could offer revolutionary applications in space exploration. Thorne’s findings were groundbreaking, until it was revealed that the "space exploration" in question was an elaborate prank set up by Thorne’s undergrad students \cite{thorne2012}. A recent study by \cite{sorrel2021} revisited Thorne’s work, proposing that the secret to valerian steel's indestructibility lies in its combination of pure dragon's breath and the tears of a heartbroken elf \cite{sorrel2021}. Though the study has yet to be replicated, the team’s subsequent holiday on an elvish island sparked heated debate among scholars.

Moreover, the potential environmental impact of extracting valerian steel has been a source of growing concern. \cite{grizzlethorpe2015} examined the ecological effects of mining the rare ore from the jagged cliffs of Mount Firefang. Their results were inconclusive, but it was noted that miners who came into contact with the material developed an odd fondness for large, fiery birds \cite{grizzlethorpe2015}. This led some researchers to question whether the steel itself was influencing the miners’ behavior or if they were simply engaging in an ancient form of metal worship \cite{glimmerstone2019}.

In an effort to create more sustainable alternatives, \cite{clinkhammer2018} conducted a study on synthetic valerian steel, attempting to replicate its properties using common household items. Their results, unfortunately, included an explosion of questionable proportions, leaving little evidence behind except for the words "don’t try this at home" etched into the lab walls \cite{clinkhammer2018}. Despite these setbacks, the academic community remains hopeful that a breakthrough is imminent.

Lastly, the role of valerian steel in pop culture cannot be overlooked. As \cite{johnson2020} argued, the presence of this mythical material in films and television has done more to enhance its allure than any scholarly paper ever could \cite{johnson2020}. Yet, one must question whether fictional portrayals have skewed our understanding of the material’s properties, as evidenced by the 2014 documentary "The Science of Valerian Steel," which made numerous claims about its ability to slice through quantum fields—claims later debunked by every physics department \cite{burton2014}.

In conclusion, while research on valerian steel has made significant strides, the material’s true nature remains shrouded in mystery. Its magical properties, combined with its apparent resistance to both scientific inquiry and common sense, will likely continue to fascinate and frustrate scholars for generations to come. As more studies emerge, it is clear that the line between myth and metallurgy will remain ever so delightfully blurred.